\SourceLeaf{67}Porro ex $\xi\xi+\eta\eta+\zeta\zeta=1$ deducimus
\[
\xi\frac{\mathrm{d}\xi}{\mathrm{d}x}+\eta\frac{\mathrm{d}\eta}{\mathrm{d}x}+\zeta\frac{\mathrm{d}\zeta}{\mathrm{d}x}=0
\]
atque hinc
\[
\frac{\mathrm{d}\frac{\eta\eta+\zeta\zeta}{\zeta}}{\mathrm{d}x}=\eta\frac{\mathrm{d}\frac{\eta}{\zeta}}{\mathrm{d}x}+\frac{\eta}{\zeta}.\frac{\mathrm{d}\eta}{\mathrm{d}x}+\frac{\mathrm{d}\zeta}{\mathrm{d}x}=\eta\frac{\mathrm{d}\frac{\eta}{\zeta}}{\mathrm{d}x}-\frac{\xi}{\zeta}.\frac{\mathrm{d}\xi}{\mathrm{d}x}
\]
Substitutis his valoribus in coëfficiente ipsius $\delta x$ in expressione pro $V$, ille fit
\[
=\xi(\frac{\mathrm{d}\xi}{\mathrm{d}x}+\frac{\mathrm{d}\eta}{\mathrm{d}y})
\]
Prorsus simili modo coëfficiens ipsius $\delta y$ in eadem expressione transit in
\[
\eta(\frac{\mathrm{d}\xi}{\mathrm{d}x}+\frac{\mathrm{d}\eta}{\mathrm{d}y})
\]
Hoc itaque pacto nanciscimur
\[
V=(\xi\delta x+\eta\delta y+\zeta\delta z)(\frac{\mathrm{d}\xi}{\mathrm{d}x}+\frac{\mathrm{d}\eta}{\mathrm{d}y})
\]

\Article{25}
Antequam ulterius progrediamur, significationem geometricam expressionum erutarum illustrare conveniet. Ad hunc finem directiones varias hic occurrentes intuitioni faciliori subiiciemus sequendo eum modum, quem in Disquiss. gen. circa superficies curvas introduximus, puta referendo illas ad puncta superficiei sphaericae radio $=1$ circa centrum arbitrarium descriptae. Primo itaque directiones axium coordinatarum $x,y,z$ denotabimus per puncta $(1)$, $(2)$, $(3)$; dein directionem normalis in superficiem et respectu spatii $s$ extrorsum ductae per punctum $(4)$; denique directionem rectae a quolibet superficiei puncto versus ipsius locum variatum ductae, per punctum $(5)$. Variationem loci ipsam, seu quantitatem $\BareRoot{(\delta x^2+\delta y^2+\delta z^2)}$, semper positive sumendam, brevitatis caussa per $\delta e$ denotabimus, arcumque inter duo sphaerae puncta, ut e. g. $(1)$ et $(5)$, sive angulum, qui illum arcum mensurat, ita $(1,5)$ scribemus. Erit itaque
\[
\delta x=\delta e.\cos(1,5),\quad \delta y=\delta e.\cos(2,5),\quad \delta z=\delta e.\cos(3,5)
\]

Haec pro quovis superficiei puncto valent. In eius limite, seu peripheria $P$, duae aliae directiones accedunt. Primo directio elementi $\mathrm{d}P$, cui respondeat
\SourceLeaf{68}punctum $(6)$; dein directio rectae huic normalis superficiem tangentis eiusque respectu introrsum ductae, cui respondeat punctum $(7)$. Per hypothesin nostram puncta $(6)$, $(7)$, $(4)$ eodem ordine iacent, ut $(1)$, $(2)$, $(3)$, observetur praeterea, $(4,6)$, $(4,7)$, $(6,7)$ exhibere quadrantes seu angulos rectos. Ita prodeunt aequationes iam supra (art. 20) traditae
\[
\begin{gathered}
\eta Z-\zeta Y=\cos(1,7),\qquad \zeta X-\xi Z=\cos(2,7),\\
\xi Y-\eta X=\cos(3,7)
\end{gathered}
\]
formulaeque art. praec. has formas induunt:
\[
\begin{aligned}
Q&=-\delta e.\cos(5,7)\\
V&=\delta e.\cos(4,5).(\frac{\mathrm{d}\xi}{\mathrm{d}x}+\frac{\mathrm{d}\eta}{\mathrm{d}y})
\end{aligned}
\]
Exprimit itaque $Q$ translationem cuiusvis puncti peripheriae $P$ a plano hanc tangente superficiei $U$ normali, in plaga ab hac aversa positive sumendam; factor ipsius $V$ autem $\delta e.\cos(4,5)$ manifesto indicat translationem cuiusvis puncti superficiei $U$ a plano hanc tangente, positive sumendam in plaga a spatio $s$ aversa.

Sed etiam factorem alterum ipsius $V$ per significationem geometricam explicare licet. Habemus enim
\[
\begin{gathered}
\xi=-\zeta.\frac{\mathrm{d}z}{\mathrm{d}x},\qquad \eta=-\zeta.\frac{\mathrm{d}z}{\mathrm{d}y}\\
\frac{1}{\zeta\zeta}=1+(\frac{\mathrm{d}z}{\mathrm{d}x})^2+(\frac{\mathrm{d}z}{\mathrm{d}y})^2
\end{gathered}
\]
Hinc prodit
\[
\begin{aligned}
\mathrm{d}\zeta&=\xi\zeta\zeta\mathrm{d}\frac{\mathrm{d}z}{\mathrm{d}x}+\eta\zeta\zeta\mathrm{d}\frac{\mathrm{d}z}{\mathrm{d}y}\\
\frac{\mathrm{d}\xi}{\mathrm{d}x}&=-\zeta\frac{\mathrm{dd}z}{\mathrm{d}x^2}-\frac{\mathrm{d}z}{\mathrm{d}x}.\frac{\mathrm{d}\zeta}{\mathrm{d}x}\\
&=-\zeta\frac{\mathrm{dd}z}{\mathrm{d}x^2}+\xi\xi\zeta\frac{\mathrm{dd}z}{\mathrm{d}x^2}+\xi\eta\zeta\frac{\mathrm{dd}z}{\mathrm{d}x.\mathrm{d}y}\\
&=-\zeta(\eta\eta+\zeta\zeta)\frac{\mathrm{dd}z}{\mathrm{d}x^2}+\xi\eta\zeta\frac{\mathrm{dd}z}{\mathrm{d}x.\mathrm{d}y}\\
\frac{\mathrm{d}\eta}{\mathrm{d}y}&=-\zeta\frac{\mathrm{dd}z}{\mathrm{d}y^2}+\eta\eta\zeta\frac{\mathrm{dd}z}{\mathrm{d}y^2}+\xi\eta\zeta\frac{\mathrm{dd}z}{\mathrm{d}x.\mathrm{d}y}\\
&=-\zeta(\xi\xi+\zeta\zeta)\frac{\mathrm{dd}z}{\mathrm{d}y^2}+\xi\eta\zeta\frac{\mathrm{dd}z}{\mathrm{d}x.\mathrm{d}y}
\end{aligned}
\]
et proin
\[
\begin{aligned}
\frac{\mathrm{d}\xi}{\mathrm{d}x}+\frac{\mathrm{d}\eta}{\mathrm{d}y}=-\zeta^3\big\{&\frac{\mathrm{dd}z}{\mathrm{d}x^2}[1+(\frac{\mathrm{d}z}{\mathrm{d}y})^2]-\frac{2\mathrm{dd}z}{\mathrm{d}x.\mathrm{d}y}.\frac{\mathrm{d}z}{\mathrm{d}x}.\frac{\mathrm{d}z}{\mathrm{d}y}\\
&+\frac{\mathrm{dd}z}{\mathrm{d}y^2}[1+(\frac{\mathrm{d}z}{\mathrm{d}x})^2]\big\}
\end{aligned}
\]
\SourceLeaf{69}cuius expressionis valorem constat esse
\[
=\frac{1}{R}+\frac{1}{R'}
\]
denotantibus $R,R'$ radios curvaturae extremos in puncto de quo agitur, et quidem positive accipiendos, quoties convexitas superficiei extrorsum vertitur.

\Article{26}
Examen attentum analysis nostrae inde ab art. 22 patefaciet suppositionem tacitam illi adhaerentem, scilicet quibusvis valoribus coordinatarum $x,y$ unicum tantummodo valorem ipsius $z$ respondere, atque valorem ipsius $\zeta$ ubique per totam superficiem $U$ esse positivum. Nihilominus veritas theorematis finalis, ad quod analysis ista perduxit, puta (I)
\[
\delta U=-\int\delta e.\cos(5,7).\mathrm{d}P+\int\delta e.\cos(4,5).(\frac{1}{R}+\frac{1}{R'})\mathrm{d}U
\]
ad hanc suppositionem non restringitur, sed generaliter valet. Quam generalitatem si statim ab initio amplecti voluissemus, vel quasdam ambages incurrere, vel methodum aliquantum diversam sequi oportuisset: sed ad eundem finem etiam per considerationes sequentes facile pervenire licet.

Analysis nostra manifesto independens est a suppositione, quod axis coordinatarum $z$ est verticalis, quin potius in illa situs axium prorsus arbitrarius manet, veritasque theorematis stabilita est pro omnibus superficiebus, pro quibus complexus omnium punctorum $(4)$ unico hemisphaerio includi potest; sufficit enim, talis hemisphaerii centrum (polum) pro $(3)$ adoptare.

Si vero proponitur superficies huic conditioni non satisfaciens, certe in duas pluresve partes dispesci poterit, quae singulae tali conditioni satisfaciant. Iam facile perspicietur, si superficies quaedam in duas partes divisa fuerit, veritatem theorematis pro figura tota statim sequi e veritate pro singulis partibus. Constet enim figura $U$ e partibus $U',U''$, sitque $P'$ peripheria figurae $U'$, atque $P''$ peripheria figurae $U''$; porro habeant $P',P''$ partem communem $P'''$, ita ut $P'$ constet ex $P'''$ et $P''''$, $P''$ vero ex $P'''$ et $P'''''$, unde manifesto peripheria figurae $U$ integra $P$ constabit ex $P''''$ et $P'''''$ Ita erit quidem
\[
\begin{aligned}
\int\delta e.\cos(5,7)\mathrm{d}P'&=\int\delta e.\cos(5,7)\mathrm{d}P'''+\int\delta e.\cos(5,7)\mathrm{d}P''''\\
\int\delta e.\cos(5,7)\mathrm{d}P''&=\int\delta e.\cos(5,7)\mathrm{d}P'''+\int\delta e.\cos(5,7)\mathrm{d}P'''''
\end{aligned}
\]
\SourceLeaf{70}sed probe notandum, valorem integralis $\int\delta e.\cos(5,7)\mathrm{d}P'''$, quatenus est pars prioris integralis, exacte oppositum esse valori eiusdem integralis, quatenus est pars posterioris integralis, quum cuivis puncto lineae $P'''$, in his duobus casibus directionibus oppositis describendae, loca puncti $(7)$ opposita adeoque valores oppositi factoris $\cos(5,7)$ respondeant. In additione itaque hae partes sese destruunt, fitque
\[
\int\delta e.\cos(5,7)\mathrm{d}P'+\int\delta e.\cos(5,7)\mathrm{d}P''=\int\delta e.\cos(5,7)\mathrm{d}P
\]
unde, quum habeatur $\delta U=\delta U'+\delta U''$, valor ipsius $\delta U$ cum formula allata (I) conspirans sponte demanat, dum haec formula cum valoribus variationum $\delta U',\delta U''$ quadrare supponitur.

Denique observamus, veritatem theorematis (I) etiam e considerationibus geometricis hauriri potuisse, et quidem facilius quam per methodum analyticam, quam tamen hic ideo praetulimus, ut occasionem, calculo variationum, pro integralibus duplicibus limitibus variabilibus inclusis parum hactenus exculto, aliquid lucis effundendi arriperemus, methodum alteram geometricam satis obviam lectori perito relinquentes.

\Article{27}
Superest, ut variationes evolvamus, quas elementa reliqua expressionis $W$ per variationem figurae spatii $s$ patiuntur, et primo de variatione voluminis spatii $s$ agemus.

Resumamus duo triangula in art. 21 considerata. iungamusque laterum puncta respondentia per rectas, ut oriatur solidum, cuius loco accipere licet prisma basis $\mathrm{d}U$, altitudinis $\xi\delta x+\eta\delta y+\zeta\delta z=\delta e.\cos(4,5)$, et quidem haec forma dabit altitudinem in forma positiva seu negativa, prout triangulum transpositum et proin totum solidum iacet extra vel intra spatium $s$. Hinc habemus (II)
\[
\delta s=\int\mathrm{d}U.\delta e.\cos(4,5)
\]
Porro hinc sequitur, variationem integralis $\int z\mathrm{d}s$ esse (III)
\[
\delta\int z\mathrm{d}s=\int z\mathrm{d}U.\delta e.\cos(4,5)
\]

Quod vero attinet ad variationem quantitatis $T$, ante omnia observamus, quum $P$ denotet limitem communem superficierum $T,U$, transpositiones puncto\SourceLeaf{71}rum peripheriae $P$ satisfacere debere huic conditioni, ut loca nova in superficie spatii $S$ maneant. Manifesto itaque per transpositionem elementi $\mathrm{d}P$, superficies $T$ patitur mutationem $\pm\mathrm{d}P.\delta e.\sin(5,6)$, perspicieturque facile, generaliter loquendo signum positivum vel negativum a signo quantitatis $\cos(4,5)$ pendere. Sed concinnius haec variatio exprimitur introducendo directionem novam, quae sit in plano superficiem spatii $S$ tangente, lineae $P$ normalis, et respectu spatii $s$ extrorsum ducta. Denotando per $(8)$ punctum huic directioni respondens, variatio superficiei $T$ a transpositione elementi $\mathrm{d}P$ oriunda erit $\mathrm{d}P.\delta e.\cos(5,8)$, sive (IV)
\[
\delta T=\int\mathrm{d}P.\delta e.\cos(5,8)
\]
ubi signum factoris $\cos(5,8)$ sponte decidet, utrum mutatio sit incrementum an decrementum.

Quum punctum $(6)$ sit polus circuli maximi per puncta $(7)$, $(8)$ ducti, punctumque $(5)$ iaceat in circulo maximo per puncta $(6)$, $(8)$ ducto, puncta $(5)$, $(7)$, $(8)$ formabunt triangulum in $(8)$ rectangulum, eritque adeo $\cos(5,7)=\cos(5,8).\cos(7,8)$: arcus $(7,8)$ autem est mensura anguli inter duo plana superficies spatiorum $s,S$ in eorum intersectione $P$ tangentia, et quidem inter eas horum planorum plagas, quae spatium vacuum includunt. Hunc angulum per $i$ denotabimus, unde $180^0-i$ erit angulus inter planorum plagas eas, quae spatium $s$ continent, formulaque nostra (V)
\[
\cos(5,7)=\cos(5,8).\cos i.
\]

\Article{28}
E combinatione formularum I....IV prodit variatio expressionis $W$
\[
\begin{aligned}
\delta W={}&\int\mathrm{d}U.\delta e.\cos(4,5).[z+\alpha\alpha(\frac{1}{R}+\frac{1}{R'})]\\
&-\int\mathrm{d}P.\delta e.\cos(5,8).(\alpha\alpha\cos i-\alpha\alpha+2\beta\beta)
\end{aligned}
\]
ubi integrale prius extendi debet per omnia elementa $\mathrm{d}U$ partis liberae superficiei spatii $s$, vel partium liberarum (si forte plures separatae adsint), integrale posterius autem per omnia elementa $\mathrm{d}P$ lineae vel linearum, quae illam partem liberam, vel illas partes liberas a reliquis spatio $S$ contiguis separant.

Iam quum in statu aequilibrii valor ipsius $W$ debeat esse minimum, adeoque admittere nequeat mutationem negativam pro ulla mutatione infi%
